% v20: complete matching parameterization, with simultaneous nonnegative high D reduced to a nonsingular sector.
% Singular positive charts remain relevant to the broader local classification.
\subsection{三高分母の Schur 補行列と，符号原理における整数性}
\label{nsch18:three-high}

ここでは $q_0=1$，$q_1,q_2,q_3>1$ とし，
$A=\operatorname{diag}(q_i)+(p_{ij})$ が
\[
 A\mathbf1=0,\qquad q_iq_j-p_{ij}p_{ji}=1
\]
を満たす局所問題を扱う．目標は，整数模型において
$D_i=-q_i-\sum_{j\ne i}p_{ji}$，$i=1,2,3$，の少なくとも一つが負になるか，
少なくとも三つ全てが狭義正となることはないかを決めることである．
本項ではいずれも一般定理とは主張しない．

\paragraph{仮定を追加しない一般帰着．}
$r_i=p_{0i}$，$x_i=p_{i0}$ とすると
\[
 \sum_{i=1}^3r_i=-1,\qquad q_i=1+r_ix_i.
\]
全ての $r_i,x_i$ は非零であり，各対 $r_i,x_i$ は同符号である．
分母一の diagonal entry を pivot とした Schur 補行列
\[
 B=A_{\{1,2,3\},\{1,2,3\}}-xr^{\mathsf T}
\]
は整数行列で，
\begin{equation}\label{nsch18:unit-diagonal}
 B_{ii}=1,\qquad B\mathbf1=0.
\end{equation}
実際，diagonal は $q_i-r_ix_i=1$ であり，row sums は
$-x_i-x_i\sum r_j=0$ である．従って，ある整数 $\alpha,\beta,\gamma$ により
\[
 B=\begin{pmatrix}
 1&\alpha&-1-\alpha\\
 \beta&1&-1-\beta\\
 \gamma&-1-\gamma&1
 \end{pmatrix}
\]
と表せる．残る三つの pair equations は
\begin{equation}\label{nsch18:linear-x}
 x_i(r_i-B_{ji}r_j)+x_j(r_j-B_{ij}r_i)=B_{ij}B_{ji}
 \qquad(1\leq i<j\leq3)
\end{equation}
であり，$B,r$ を固定すれば $x$ に関する線形方程式である．
これは $p_{ij}=B_{ij}+x_ir_j$ を元の pair equation へ代入して得られる．
また，column sums を取ると
\begin{equation}\label{nsch18:defect}
 D_0=-1-\sum_i x_i,\qquad
 D_i=-\sum_{j=1}^3B_{ji}+r_iD_0.
\end{equation}
従って $D_0>0$ なら $\sum_{i=1}^3D_i=-D_0<0$ であり，
高分母側の少なくとも一つは直ちに負となる．
さらに，正の weights を持つ場合は全 $D$ ベクトルが零である整数模型は存在せず，
$D_0=0$ でも高側に負の係数が現れる．
従って正の weights の下で残る場合は $D_0<0$ に限られる（\ref{nhd20:reduction}）．

この行列問題は，二次の principal minors が全て一である非対称行列の特殊な場合である．
Lin--Sturmfels\cite[\S2, Proposition~4]{LinSturmfels} は principal minors と cycle sums の
整数多項式による変換を与える．本問題ではさらに row sums，整数性，正の weights を課すため，
一般の principal-minor relations だけから係数の符号や整除性は結論しない．

\paragraph{全ての正整数三高分母模型を尽くす canonical 表示．}
完全 matching 定理（\ref{nm18:matching}）により，正の numerators の辺は
高 rays 間の $ab$ と，残る高 ray $c$ と低 ray $0$ を結ぶ $c0$ に限られる．
従って以下の表示は追加の matching 仮定を必要とせず，全ての正整数三高分母模型を扱う．
これは有限個の模型への分類ではなく，残る整数条件を共通の線形系に表す帰着である．

$s,t,S,T,U\ge1$，$u=s+t-1$ と置く．適切に $a,b$ を交換すれば $k\ge0$ を選べ，
\[
 (q_a,q_b,q_c,q_0)=(1+sS,1+tT,1+uU,1),
\]
\[
 P=\begin{pmatrix}
 0&tS+\alpha&-(uS+\alpha+1)&-S\\
 sT+\beta&0&-(uT+\beta+1)&-T\\
 -(sU-k)&-(tU+k+1)&0&U\\
 -s&-t&u&0
 \end{pmatrix}
\]
となる．$H=tS+\alpha,K=sT+\beta,L=uS+\alpha+1,N=uT+\beta+1,
M=sU-k,O=tU+k+1$ の六つが全て正であることを保持する．
$\alpha,\beta$ は整数であり，
$\mathcal A=s\alpha-t,\mathcal B=t\beta-s,\mathcal C=s+uk$ と置く．
残る三つの pair equations は
\[
 \begin{pmatrix}
  \mathcal B&\mathcal A&0\\
  -\mathcal C&0&\mathcal A+1\\
  0&\mathcal C-1&\mathcal B+1
 \end{pmatrix}
 \begin{pmatrix}S\\T\\U\end{pmatrix}
 =\begin{pmatrix}
 -\alpha\beta\\k(\alpha+1)\\-(\beta+1)(k+1)
 \end{pmatrix}
\]
にちょうど同値である．正の weights は
$w=(q_a,H+1/(2K),u+1/(2U),1)$ と明示できる．
幾何的な奇数分母条件は $sS,tT,uU$ が全て偶数であることとして保持する．
公開算術稿の \texttt{sh14:matching-setup} と \texttt{sh14:linear-reduction} は，
この表示と線形化の証明箇所である．matching 定理により全正整数三高模型へ適用できる．

この canonical 表示では
\[
 \mathcal A=s\alpha-t\geq1,\qquad
 \mathcal B=t\beta-s\leq-2,\qquad
 \mathcal C=s+uk\geq2,\qquad \beta\leq-1.
\]
$d=U+1-S-T=-D_0$ と置くと
\[
 D_a=sd-\beta-k-1,\quad D_b=td-\alpha+k,\quad
 D_c=\alpha+\beta+1-ud.
\]
従って全ての高係数が非負となるには
\[
 sd\geq\beta+k+1,\quad td\geq\alpha-k,\quad
 ud\leq\alpha+\beta+1
\]
が必要十分である．全て狭義正なら右辺をそれぞれ一だけ強める．
境界 $\beta=-1$ では $k=0,D_0>0$ が証明され，その範囲から全高係数非負の模型は生じない．
全ての正整数三高分母模型のうち，高係数符号の問題で今後調べる範囲は
$b=-\beta\ge2$，$d\ge1$，$\alpha\ge k\ge b+1$，$\Delta>0$ である．
全ての特異系はこの符号問題から除かれる（\ref{nhd20:reduction}）．
matching 性が解決したことと，この $D$ 符号の問題が解決したことを混同しない．

固定 $s\geq2,t\geq1$ に対する $\beta=-1$ の全解は，
\[
 K_0=s^2+st+t^2-s-t,\quad g=\gcd(s-1,t),\quad G\mid K_0g,\quad G>0
\]
を列挙し，$\alpha=(G+K_0)/(st)$ の正整数性と
\[
 S=\frac{\alpha(s-1)(s\alpha-t+1)}G,\quad
 T=\frac{\alpha s(s+t-1)}G,\quad
 U=\frac{\alpha s(s-1)}G
\]
の正整数性を課して得る．幾何的な奇数条件は $sS,tT,uU$ が全て偶数であることとして最後に課す．
この完全な有限手続きは，固定 $s,t$ に関するものであり，一般分母の有限性は主張しない．

\paragraph{特異系は実在する：分類から除外しない．}
係数行列の determinant は
\[
 \Delta=-\mathcal B(\mathcal A+1)(\mathcal C-1)
          +\mathcal A\mathcal C(\mathcal B+1).
\]
$\Delta\ne0$ では Cramer 分子の整除性を確認する．$\Delta=0$ では拡大行列のrankと整数解格子を調べる．
実際，$(s,t,\alpha,\beta,k)=(3,12,12,-4,1)$ では rank が二で，正整数全解は
\[
 (S,T,U)=(136+200n,291+425n,93+136n),\qquad n\geq0
\]
で尽くされる．分母は $(409+600n,3493+5100n,1303+1904n,1)$，
高係数は $(-997-1467n,-4007-5868n,4671+6846n)$ である．
この固定sliceは解決済みであり，同じ補助データの探索を繰り返さない．
全ての特異系の分類と，大域の五 core 接着は未解決である．
ただし特異系の高係数が全て非負になる可能性は排除済みであり，その候補として再探索しない．
正整数三高分母の非matching型は存在しないことが証明されており，再探索の対象ではない．

\paragraph{高係数の同時非負性は非特異領域だけに残る．}
\label{nhd20:reduction}
canonical 表示で $b=-\beta$，$d=-D_0=U+1-S-T$ と置く．
整数 row sums，pair products，正の weights のみを仮定すると，
\[
 D_a,D_b,D_c\ge0
 \quad\Longrightarrow\quad
 b\ge2,\quad d\ge1,\quad \alpha\ge k\ge b+1,
\]
\[
 \Delta=\mathcal C(\mathcal A-\mathcal B)
          +\mathcal B(\mathcal A+1)
 \ge(-\mathcal B)(-\mathcal B-1)>0.
\]
分母の上限，奇数性，clutching，node 条件は不要である．
したがって，階数二を含む全ての特異 Schur 系から
全高 $D$ 非負の模型が生じる可能性はない．
この結論は，特異な正整数局所模型そのものを排除するものではない．

再利用する証明の要点を残す．
\[
 D_a=sd+b-k-1,\quad
 D_b=td-\alpha+k,\quad
 D_c=\alpha-b+1-ud,\qquad
 D_a+D_b+D_c=d
\]
である．まず全 $D$ 零を除く．その場合は
$d=0,\alpha=k=b-1$ で，$b\ge2$ である．
\[
 m=b^2-b+1,\qquad K=s^2+st+t^2-s-t,\qquad
 r=(-s,-t,u)^{\mathsf T},\quad x=(-S,-T,U)^{\mathsf T},
\]
\[
 B_0=\begin{pmatrix}
 1&b-1&-b\\
 -b&1&b-1\\
 b-1&-b&1
 \end{pmatrix}
\]
と置くと，Cramer 恒等式は $\Delta=mK$，$K=m-1$，
$x=r-B_0r/m$ を与える．整数性から各成分で $m\mid B_0r$ である一方，
\[
 B_0^{\mathsf T}B_0=m(3I-J),\qquad
 \lVert r\rVert^2=2m-1,\qquad \mathbf1^{\mathsf T}r=-1
\]
である．ここで $J$ は全成分が一の三次行列である．従って
$m^2\mid\lVert B_0r\rVert^2=m(6m-4)$，従って $m\mid4$ となる．
$m$ は三以上の奇整数だから矛盾する．
特に $D_0=0$ なら高係数には正と負がともに現れる．

同時非負性には $d\ge1$ が必要で，
$0\le D_a\le d$ から
\[
 (s-1)d+b-1\le k\le sd+b-1
\]
を得る．$k=b-1$ なら $s=1$，$u=t$ となり，
線形系の第三式は
\[
 t(b-1)T-tbU=b(b-1)>0
\]
となる．しかし $U=S+T+d-1>T$ なので左辺は負である．
$k=b$ なら $(s-1)d\le1$ であり，$s=1$ または $(s,d)=(2,1)$ となる．
前者では第三式 $tb(T-U)=b^2-1>0$ が同じ理由で矛盾する．
後者では $\alpha=b+t$，$U=S+T$ であり，第一式と第三式の線形結合が
\[
 t(1+b^2+tb)S=t+2b-b^3<t
\]
を要求するため，$S\ge1$ に反する．
この境界排除は determinant で割らないので，途中で特異系を落としていない．

以上から $k\ge b+1$ であり，
$D_c\ge0$ と $k\le sd+b-1$ から
$\alpha-k\ge(t-1)d\ge0$ が従う．
$\bar b=-\mathcal B=s+tb$ と置けば
\[
 \mathcal C-\bar b=(s-1)b+u(k-b)\ge0,\qquad
 \Delta=\bar b(\bar b-1)+(\mathcal C-\bar b)(\mathcal A+\bar b)
\]
により determinant の狭義正値性を得る．

残る非特異領域での一般符号問題は未解決である．
この領域の整数局所模型自体は存在する．例えば
\[
 (s,t,\alpha,\beta,k)=(1,2,10,-2,5),\qquad
 (S,T,U)=(4,5,11)
\]
は $q=(5,11,23,1)$，$D=(-1,1,3,-3)$，
正の weights $(5,109/6,45/22,1)$ を与える．
従って $\alpha\ge k\ge b+1$ と $\Delta>0$ だけで局所模型を捨てず，
Cramer 整除性と全高 $D$ の非負性を引き続き課す必要がある．
\nolinkurl{verify_high_d_boundary_v20.py} は
四十二の多項式恒等式と三つの厳密行列回帰を検算する．
一般命題の根拠は上の恒等式・不等式であり，有限範囲の代入結果ではない．

\paragraph{五つの整数を固定した三高 $D$ の完全復元．}
\label{nhd21:reconstruction}
同時非負性の残余問題では，$x=D_a,y=D_b,z=D_c\ge0$ を直接指定する．
$s,t$ も固定することが必要であり，$D$ だけを固定した有限性とは主張しない．
次の変数は再開後の復元と消去にそのまま使える：
\[
 u=s+t-1,\quad d=x+y+z\ge1,\quad
 k=sd+b-1-x,\quad\alpha=ud+b-1+z,
\]
\[
 h=s-x,\quad j=t-y,\quad\ell=u+z,\quad
 \lambda=ux+sz,\quad\mu=uy+tz,
\]
\[
 F=b-1+\lambda,\quad m=F^2+F+1,\quad
 E=\mu h-\ell-j,\quad J=h+\ell-\lambda j.
\]
正の模型では $h\ge1$，すなわち $D_a\le s-1$ が新しく必要となる．
$j$ の符号は制限しない．$E,J$ は $b$ に依存せず，
\[
 S=h+\frac{FE-J}{m},\quad
 T=j+\frac{E+(F+1)J}{m},\quad
 U=\ell+\frac{(F+1)E+FJ}{m}
\]
が全尺度を一意に復元する．三尺度の整数性は単一条件 $m\mid FE-J$ と同値である．
元の三つの高高 pair 残差はこの尺度で互いに等しくなるが，自動的に零にはならない．
\[
 (-s-tb)S+(s\alpha-t)T-\alpha b=0
\]
を最後まで保持し，全尺度の正値性，canonical 六符号量の正値性，
$\alpha\ge k\ge b+1$ も確認する．これらは元の模型への必要十分判定となり，
正の weights は公開稿の明示式から得られる．

\paragraph{ノルムによる全 $b$ の有限化．}
$v=S-h,w=T-j$ と置くと
\[
 E=(F+1)v+w,\qquad J=-v+Fw,
\]
\[
 N=E^2+EJ+J^2=m(v^2+vw+w^2).
\]
零ノルムは，正尺度の下では
$(s,t,x,y,z)=(3,3,0,0,1),(2,2,1,1,1)$ に限られる．
元の残余式がそれぞれ $b(b+5)=27$，$b(b+9)=14$ を要求するため，両方とも排除される．
従って $N>0$ であり，
\[
 m\mid N,\qquad m\le N,\qquad
 b\le 1-\lambda+
 \left\lfloor\frac{\sqrt{4N-3}-1}{2}\right\rfloor
\]
となる．実装では平方根を \texttt{isqrt} で厳密に評価する．
$s,t,x,y,z$ を固定すれば，任意 cutoff なしに $b\ge2$ を全て調べられる．
この五つの整数そのものへの一様上界は未証明である．

\paragraph{誤った簡約を防ぐ二つの反例．}
$(s,t,x,y,z,b)=(1,1,0,3,0,3)$ では
$F=2,m=7,E=4,J=2,N=28$ である．$m\mid N$ でも
$m\nmid FE-J=6$ なので尺度は整数でない．ノルム整除性だけに置き換えない．
また $(1,3,0,4,0,2)$ では $(S,T,U)=(3,5,11)$ は正整数で，
六符号量も正，$D=(0,4,0,-4)$ であるが，元の三つの高高 pair 残差は全て三である．
この例は模型ではなく，残余式を落とせないことの確認例である．

\paragraph{再利用する実装と次の課題．}
\nolinkurl{verification/verify_high_d_reconstruction_v21.py} の
\texttt{reconstruct\_fixed\_data(s,t,x,y,z)} が全 $b$ の完全復元器である．
通常の実行は二十四の多項式係数恒等式と五つの厳密回帰を確認する．
\texttt{--fixed 3 4 0 1 1} のように固定データを渡せる．
既存の有限族検査は $s,t\le12,0\le x<s,0\le y,z\le4,d\ge1$ の
$23256$ 組について $b\ge2$ を全範囲で尽くし，模型は零だった．
結果と実行当時の source は \nolinkurl{research/} に保存した．
これは五変数に制限を付けた結果であり，一般不存在ではない．
第22版ではこの復元式に entry 正値性を加え，固定するデータを $(s,t)$ の二つへ減らす．
次の節の有限化と合わせて用い，二変数自体が動く問題を残す．

% v22 research notes: retain only material useful for future work.
\subsection{三高同時非負性：低行二変数までの有限化}
\label{notes:hd22}

\paragraph{現在の到達点}

第21版の五固定データ $(s,t,x,y,z)$ ごとの $b$ 有限化は，
第22版で固定 $(s,t)$ ごとの全 $x,y,z,b$ の有限復元へ強まった．
新しい主モジュールは
\nolinkurl{high_d_fixed_low_row_finiteness_v22.tex} である．
低行 $(-s,-t,s+t-1,0)$ の二整数自体の一様上限，
および三高同時非負模型の一般不存在は未解決である．

利用する記号と核となる式だけを残す．
\[
 u=s+t-1,\ h=s-x\ge1,\ B=b-1\ge1,\
 F=B+ux+sz,\ m=F^2+F+1.
\]
固定 $(s,t,x,z,b)$ において
\[
 a=uh-B(F+1),\qquad
 n_0=(F+1)(Bt+h)+F(u+z)+htz,
 \qquad mT=n_0+ay.
\]
$K=sT-b>0$ と $y\ge0$ が
\[
 \boxed{F\le M:=st+(2u+1)h-2}
\]
を与える．上限の証明は $T$ の正整数性を過剰に使わず，entry $K$ の正値性を
直接使う．$F$ は整数なので，$F\ge M+1$ の排除から表示した上限が従う．
さらに $b-1+sz\le M-ux$ より
\[
 (3u+1)x\le st+(2u+1)s-3
\]
も従う．これは $x\le s-1$ の追加の絞り込みとして使える．

残余式を $y$ の多項式にして
\[
 \mathcal R(Y)=suaY^2+(aL-C)Y+C_0
\]
とすると，$L,C>0$ であり，$a=0$ でも一次係数が $-C<0$ となる．
従って各 $(s,t,x,z,b)$ に非負整数根は高々二個，$a\le0$ なら高々一個である．
二次式の係数と $C_0$ の計算法は公開用証明に明示した．
根は判別式の整数平方根と正確な割り算で全て求められ，$y$ の探索上限は不要である．

\paragraph{完全復元で残す検査}

二次式の根であることは元の第一残余そのものの零を意味する．
ただしそれだけでは整数模型ではない．必ず
\[
 d=x+y+z\ge1,\quad \alpha\ge k\ge b+1,\quad m\mid FE-J,
\]
三尺度の正値性，六つの符号量の正値性を検査する．
最後に元の $q,P$ を組み立て，row sums，六つの pair products，
$D=(x,y,z,-d)$，正 weights に対する十二の gaps を再計算すれば，
導出式の取り違えも検出できる．$j=t-y$ の非負性は仮定しない．

第21版の零ノルム分類は新手順の例外処理として不要になったが，
$(s,t,x,y,z)=(3,3,0,0,1)$ と $(2,2,1,1,1)$ の両族は
元残余を捨ててはいけないことを示す回帰例として残した．
前者の尺度 $(3,3,6)$ は $b$ によらず整数だが，
残余 $27-b(b+5)$ は $b\ge2$ の整数で零にならない．

\paragraph{再現可能な厳密計算と次の分岐}

標準ライブラリのみの
\nolinkurl{verify_high_d_low_row_finiteness_v22.py}
は二十四の多項式恒等式を係数ごとに確認する．
\texttt{--fixed s t} はその低行の全模型を完全復元し，
\texttt{--scan 12} は $1\le s,t\le12$ の全 $144$ 低行に同じ処理を行う．
探索箱がかかるのは $s,t$ だけであり，$D$ や分母は全範囲である．
結果は \nolinkurl{high_d_low_row_finiteness_verification_v22.json} にある．
全五組候補 $2{,}004{,}015$，$a=0$ の一次式 $219$，
非負整数残余根六個で，全て $d=0$ のため対象外だった．
零 $D$ の六根の $(s,t,b)$ は
\[
 (3,4,6),(4,3,6),(6,9,13),(9,6,13),(10,12,19),(12,10,19).
\]
これらの尺度は有理数であり，整数の零 $D$ 模型の反例ではない．

次の有力課題は，無上界の $(s,t)$ に対して残余二次式の非負整数根と
単一合同条件を同時に扱うことである．有限復元を使えば，
単純な分母・高係数 cutoff の拡大を避け，低行の構造に沿った反例探索ができる．
また観測範囲では $d\ge1$ の整数残余根が合同条件を課す前から存在しなかった．
これは，尺度の整数性を外した有理復元段階にもより強い障害がある可能性を示すが，
現時点では観測であって定理ではない．
判別式は
\[
 \mathcal D_y=(aL-C)^2-4suaC_0
 =C^2+a\bigl(aL^2-2LC-4suC_0\bigr)
\]
である．一般の全展開では扱いやすい積を得ていない．
次は $a$ の符号と $a=0$ の退化面を先に分け，
整数平方間の評価と合同条件を組み合わせる方がよい．
$a=0$ は $B(F+1)=uh$，すなわち
\[
 u(s-bx)=B(b+sz)
\]
を要求するので，$s-bx>0$ という追加制約も得られる．
この面で残余は $-CY+C_0$ であり，$y=C_0/C$ の整数性へ直ちに帰着する．

判別式とは別に，$s=t$ かつ $x=y$ の対称 sector は無上界で排除できた．
このとき $J=-E$，単一合同条件から $n=E/m\in\mathbb Z$，
$S=h+(F+1)n$，$T=h-Fn$ となる．
$F\ge h+1$ により正値性は $n=0$ を強制し，既知の零ノルム排除に矛盾する．
この証明は整数性と正値性の両方を使うので，
一般 sector でも $E+J=\ell(sy-tx)+h-j$ の小さな値から類似の格子拘束を
引き出せるかを検討する価値がある．
一般不存在を目標に固定せず，低行に沿った無限族・合同条件の反例・
残余式の因数分解が現れた場合は，その構造の記述へ研究計画を切り替える．

% v23 notes: keep only reusable formulas, exact scope, and useful next steps.
\subsection{三高同時非負性：非対称量と退化面}
\label{notes:hd23}

\paragraph{確定した四領域と残る量化}

主証明は \nolinkurl{high_d_asymmetry_lattice_obstruction_v23.tex} にある．
第22版の固定低行 $(s,t)$ ごとの全係数有限復元を保ち，
\[
 \delta=sy-tx,\qquad a=uh-(b-1)(F+1)
\]
を使って，次の四領域を新たに無上界で除外した：
\[
 \delta=0,\ s\ge t;\qquad
 \delta=-1;\qquad
 \delta=1,\ x\ge1;\qquad
 a=0,\ s\le2.
\]
$s=t$ だけで模型を除外したわけではない．
$\delta=0,s<t$，$|\delta|\ge2$，$\delta=1,x=0$，
一般の $a=0,s\ge3$ は残る．三高同時非負性の一般不存在は未解決である．

再利用する核は，$n=\ell-U\in\mathbb Z$ に対する
\[
 S=j-(F+1)n-\ell\delta,\qquad
 T=h+Fn+\ell\delta,
 \qquad F+\delta=b-1+h+s(d-1)
\]
である．$n$ は node 配分ではなく尺度の格子変数である．
$\delta=-1$ では正値性が $x=1,n=1,t=sy+1$ を強制し，
残る差方程式が正係数表示を持つ．
$\delta=1,x\ge1$ では $x=1,n=-1,z=0,t=sy-1$ を強制する．
後者では差方程式 $Q_+=0$ に加えて元残余を残すことにより
\[
 y\mathfrak r=\bigl(sy(y+1)+1\bigr)Q_++N_+,
 \qquad N_+>0
\]
を得る．この二方向の違いは今後の $|\delta|\ge2$ でも保持すべきである．

\paragraph{差方程式だけでは通過する新しい回帰例}

\[
 (s,t,x,y,z,b)=(5,4,1,1,0,3)
\]
は $\delta=1$ であり，復元尺度は $(S,T,U)=(6,2,9)$，
$\alpha=18,k=11$，$q=(31,9,73,1)$ となる．
単一合同と六つの符号条件を満たすが，三つの元の高--高 pair 残余は全て $16$ である．
これは模型ではない．単一合同と正値性に元残余を付け加える必要を具体的に示す．
また $(s,t,x,y,z,b)=(3,6,1,2,0,3)$ では
$\delta=0,s<t$，$(S,T,U)=(4,2,8)$，元残余 $-18$ である．
こちらも模型ではなく，今回の比例領域の順序仮定を勝手に捨てる根拠にならない．

\paragraph{退化面の有界整数と追加の整除条件}

$a=0$ では
\[
 s\ell=bF,\qquad uh=(b-1)(F+1),\qquad
 u(s-bx)=(b-1)(b+sz)>0
\]
となる．公開証明の補題により
\[
 K_0=h(s-1)+1,\qquad R_0=b^2-sT
\]
は
\begin{equation}\label{notes:hd23-degenerate-quadratic}
 \begin{gathered}
 1\le R_0\le K_0,\\
 R_0F^2+(R_0-bK_0+sh)F+R_0-b^2+sh=0
 \end{gathered}
\end{equation}
を満たす．$s=1,2$ では $R_0$ はそれぞれ一通りと三通りだけとなり，
各判別式の法三，五，七の平方障害で全て除外できた．

元残余も使うと，さらに
\begin{equation}\label{notes:hd23-extra-divisibility}
 \boxed{\ b-1\mid h(s-1)(R_0-1)\ }
\end{equation}
を得る．この式の短い導出を残す．$B=b-1$ とし，
\[
 S_0=FT-Bt-\ell,\quad \alpha_0=F+tz,\quad
 H=s+tb,\quad K=sT-b,\quad \rho=BH-uK
\]
と置く．$S=S_0+By$，$\alpha=\alpha_0+uy$ だから
\[
 \mathfrak r=\mathfrak r_0-\rho y,
 \quad
 \mathfrak r_0=-HS_0+(s\alpha_0-t)T-\alpha_0b.
\]
$uh=B(F+1)$ より $B\mid h\rho$ なので，$\mathfrak r=0$ なら
$B\mid h\mathfrak r_0$ である．法 $B$ で $b=1$ とすると，恒等式
\[
 h(\mathfrak r_0+tT)
 =h(tT+1)(s\ell-F)+t(1-sT)uh
 \pmod B
\]
の右辺は零である．従って $B\mid htT$ であり，
$ht\equiv h(1-s)\pmod B$ から $B\mid h(s-1)T$ となる．
最後に $sT=b^2-R_0\equiv1-R_0\pmod B$ を使えば
式~\eqref{notes:hd23-extra-divisibility} が従う．

これにより，退化面の大部分は低行の一整数 $s$ だけを固定して有限化できる．
$x\ge1$ の場合は $s-bx\ge1$ から
\[
 2\le b\le\left\lfloor\frac{s-1}{x}\right\rfloor.
\]
$x=0$ かつ $R_0\ge2$ では $h=s$ なので
\[
 2\le b\le1+s(s-1)(R_0-1)\le1+s^2(s-1)^2.
\]
$s,h,b,R_0$ を固定すれば，式~\eqref{notes:hd23-degenerate-quadratic} の
正の二次係数 $R_0$ により整数 $F$ は高々二個であり，
\[
 u=\frac{B(F+1)}h,\qquad \ell=\frac{bF}s,\qquad
 t=u-s+1,\quad x=s-h,\quad z=\ell-u
\]
で残りを復元できる．$y$ は第22版の恒等零でない一次残余で決まる．
全ての整数性，正値性，単一合同，元の pair 残余は引き続き確認する．

この有限化から外れる唯一の組合せは $x=0,R_0=1$ である．
この枝には必要条件
\begin{equation}\label{notes:hd23-exceptional-congruence}
 s^2\mid2(b-1)
\end{equation}
がある．実際，$s\mid b(b-1)$ と $sT=b^2-1$ を比較すると $s\mid b-1$ であり，
$b-1=s\kappa$，$F=sL$ と書ける．
式~\eqref{notes:hd23-degenerate-quadratic} は
\[
 s\bigl(L^2+(1-\kappa K_0)L-\kappa^2+1\bigr)=2\kappa
\]
を与える．従って $s\mid2\kappa$，すなわち式~\eqref{notes:hd23-exceptional-congruence} が従う．
この枝の一般不存在は証明していない．平方判別式と元残余を併用する，
より小さい次の標的として扱える．

$a=0$ だから全ての整数尺度が存在しない，という強化は誤りである．
\[
 (s,t,x,y,z,b)=(3,138,0,0,4,16)
\]
では $F=27,m=757$，$(S,T,U)=(54,84,141)$，$R_0=4$ となり，
単一合同と六符号は成立する．元残余は $5658$ であり模型ではない．
この例は一般の退化面で元残余を使う必要性を示す．

\paragraph{残る $\delta=1,x=0$ の二整数による完全復元}

この枝では $s=y=1$ である．$k=b+z\ge b+1$ より $z\ge1$ であり，
格子正値性から $p=-n\ge1$ と書ける．
\[
 F=b-1+z,\quad m=F^2+F+1,\quad
 S=(F+1)p-z-1,\quad T=t+z+1-Fp,\quad U=t+z+p.
\]
差方程式は
\[
 (t-1)(z-1)=mp-F(z+2)
\]
である．$z=1$ なら $mp=3F$ だが，$F=b\ge2$ で
$m-3F=(F-1)^2>0$ なので不可能である．従って未解決枝は $z\ge2$ に限られる．

固定 $b,z\ge2$ に対して $D=z-1$ とし，整数 $p$ の一次式を
\[
 \tau=D+mp-F(z+2),\quad
 \sigma=(F+1)p-z-1,\quad
 \theta=\tau+D(z+1-Fp)
\]
と定める．$t=\tau/D$，$S=\sigma$，$T=\theta/D$ であり，
元残余の分母を払った式は
\begin{equation}\label{notes:hd23-axis-quadratic}
 \begin{aligned}
 \mathcal P(p)={}&-D(D+b\tau)\sigma
 +(z\tau+FD)\theta\\
 &-bD\bigl((z+1)\tau+FD\bigr)=0.
 \end{aligned}
\end{equation}
その二次係数は
\[
 (b^2+bz+z^2)m>0
\]
なので，各固定 $(b,z)$ に正整数根 $p$ は高々二個である．
整数平方根による全根計算の後，
\[
 D\mid\tau\quad\Longleftrightarrow\quad
 z-1\mid(b^2+b+1)p-3b
\]
と $t,S,T,U$ および六符号の正値性を検査すればよい．
ここでは差方程式と元残余の双方を保持しており，通過したものは
第21版の完全復元により全 pair products と正 weights を持つ．
尺度整数性の単一合同も直接再確認できる．

標準 Python 検証器の \texttt{remaining\_axis\_fixed\_b\_z(b,z)} はこの完全復元を実装する．
\texttt{--axis b z} でも呼び出せる．$p$ や $t$ に探索上限はない．
固定している $b,z$ 自体には上限を与えていないので，この手順だけで残余枝全体を
不存在と結論してはいけない．

\paragraph{検算と今後の焦点}

\nolinkurl{verify_high_d_asymmetry_v23.py} は標準ライブラリのみを用い，
四十一の多項式恒等式を係数ごとに検査する．
退化面の四つの判別式は，法 $9,25,9,49$ の全九十二剰余類で
平方剰余と交わらないことを独立確認する．七組の回帰例は
元の rows，pair products，$D$，必要な場合の正 weights を再計算する．
結果は \nolinkurl{high_d_asymmetry_verification_v23.json} にあり，
source と公開証明の hash，検算の量化，残る範囲を記録している．
第22版の二百万余の復元計算は再実行していない．

次は，$\delta$ の絶対値を単に切って列挙するのでなく，
格子正値性が $x,n$ をどう制限するかを一般の $\delta$ で整理する．
とくに負側 $\delta=-k$，$s\ge k$ では，$j>0,n\ge1$ と $S>0$ から
\[
 u(x-k)+(s-k)z<t-(b-1)-1
\]
を得るので，$x\le k$ が必要である．
正側でも $\delta=k$，$s\ge k$ なら $n\le-1$ と $T>0$ が $x\le k$ を与える．
この一般の狭い格子幅と，二次残余の符号証明を組み合わせることが有望である．
また，退化面の $x=0,R_0=1$ と，正側零境界の式~\eqref{notes:hd23-axis-quadratic} は
独立した二つの具体的な未解決問題として追える．
どちらも反例や無限族が現れれば，その構造の記述へ計画を変更する．


\paragraph{整数性を外すと，三つの高 $D$ が全て狭義正になる．}
次は\emph{有理数による代数的緩和}の反例であり，整数 branch denominator を持つ
局所 chart の反例ではない．頂点順を $(a,b,c,0)$ とし，
\[
 q=\left(\frac{79}{25},\frac{128}{25},\frac{274}{25},1\right),\qquad
 P=\begin{pmatrix}
 0&179/25&-231/25&-27/25\\
 53/25&0&-259/50&-103/50\\
 -91/25&-266/25&0&83/25\\
 -2&-2&3&0
 \end{pmatrix}
\]
と取る．全 row sums と pair equations を満たし，
\[
 w=(1,7/3,13/4,1)
\]
について全ての $w_j-w_ip_{ij}/q_i$ は狭義正である．
実際，正の辺は $a\leftrightarrow b$ と $c\leftrightarrow0$ のみで，
対応する四つの差は
\[
 \frac{16}{237},\quad\frac{13}{384},\quad
 \frac{17}{1096},\quad\frac14
\]
である．一方，
\[
 D=\left(\frac9{25},\frac9{25},\frac{23}{50},-\frac{59}{50}\right).
\]
この例でも low row と Schur 補行列
\[
 B=\begin{pmatrix}1&5&-6\\-2&1&1\\3&-4&1\end{pmatrix}
\]
は整数である．従って，$B$ と low row の整数性まで仮定しても，
$q$ と逆向き numerators の整数性を外すと一般符号原理は成立しない．
正の weights と主小行列の不等式だけによる証明は不可能であり，
残る整除性を本質的に用いる必要がある．
この反例と上記恒等式は \texttt{verify\_three\_high\_schur\_v12.py} で有理数演算により再確認できる．

\paragraph{$\beta=-1$ を持つ無限族と，一つの slice の完全な記述．}
\label{nsch18:beta-minus-one}
上の完全約数表示の $t=1$ の場合には，さらに簡単な閉じた表示がある．
一般のmatching型に $\beta\leq-2$ を仮定してはいけないことも，この族から分かる．
実際，$t=1$，$\beta=-1$，$k=0$ と固定した完全マッチング型は，
次の二つの正整数 $d,e$，$de\geq2$ でちょうど記述される：
\begin{align*}
 s&=de,&\alpha&=d(e+1),\\
 S&=d(e+1)^2(de-1),&
 T&=de(e+1),&
 U&=(e+1)(de-1).
\end{align*}
分母と numerators は
\[
 (q_a,q_b,q_c,q_0)=(1+sS,1+T,1+sU,1),
 \qquad
 P=\begin{pmatrix}
 0&S+\alpha&-(sS+\alpha+1)&-S\\
 sT-1&0&-sT&-T\\
 -sU&-(U+1)&0&U\\
 -s&-1&s&0
 \end{pmatrix}
\]
である．三つの高分母は全て奇数で，一より大きい．
$H=S+\alpha$，$K=sT-1$ と書くと，正の weights は例えば
\[
 w=\left(q_a,\ H+\frac1{2K},\ s+\frac1{2U},\ 1\right)
\]
で与えられる．正の四つの directed edges に対応する basin gaps は
\[
 \frac1{2K},\quad\frac1{2q_b},\quad\frac1{2U},\quad\frac1{2q_c}
\]
であり，残る edges は負の numerators を持つので正値性は自明である．
また
\[
 D_0=S+T-U-1=(\alpha-1)U+T-1>0,
 \qquad
 (D_a,D_b,D_c)=(-sD_0,-D_0-\alpha,sD_0+\alpha).
\]
従って，これは三つの高係数が全て正となる反例ではない．

完全性の確認も短い．$t=1,\beta=-1,k=0$ の線形系では
\[
 \Delta=s(s-\alpha),\qquad
 S=\frac{\alpha^2(s-1)}{\alpha-s},\quad
 T=\frac{\alpha s}{\alpha-s},\quad
 U=\frac{\alpha(s-1)}{\alpha-s}.
\]
$s\geq2$ であり，$\alpha=s$ の場合には $S$ の Cramer numerator
$-\alpha^2s(s-1)$ が非零なので系は不整合である．
$S>0$ には $d=\alpha-s>0$ が必要である．
$T,U$ の整数性は $d\mid s^2$ および $d\mid s(s-1)$ と同値なので，$d\mid s$ が従う．
逆に $d\mid s$ なら三つ全てが整数となる．$s=de$ と書けば表示した族を得る．
pair equations と row sums はそのまま代入して確認できる．

小さい具体例として $d=1,e=2$ を取ると
\[
 q=(19,7,7,1),\qquad
 P=\begin{pmatrix}
 0&12&-22&-9\\
 11&0&-12&-6\\
 -6&-4&0&3\\
 -2&-1&2&0
 \end{pmatrix},\qquad
 D=(-22,-14,25,11),
\]
かつ $w=(19,265/22,13/6,1)$ となる．
これは既存の有限 catalog にも含まれる模型であり，この slice 全体の無限パラメータ表示と，残る整除性の特徴づけを再利用できる．
これらの恒等式と具体例は \texttt{verify\_three\_high\_bounds\_v13.py} で厳密演算により検算する．
